\documentclass[11pt,a4paper]{article}
\usepackage[margin=22mm]{geometry}
\usepackage{amsmath,amssymb}
\usepackage{fontspec}
\usepackage{polyglossia}
\setdefaultlanguage{russian}
\setmainfont{DejaVu Serif}
\usepackage[unicode,hidelinks]{hyperref}

\allowdisplaybreaks[1]
\raggedright
\setlength{\parskip}{0.35em}
\newcommand{\dd}{\mathop{}\!\mathrm{d}}
\newcommand{\avg}[1]{\langle #1\rangle}
\newcommand{\ket}[1]{\lvert #1\rangle}
\newcommand{\bra}[1]{\langle #1\rvert}
\newcommand{\braket}[2]{\langle #1\mid #2\rangle}
\DeclareMathOperator{\Var}{Var}

\newcommand{\solutiontitle}[4]{%
  \hypersetup{pdftitle={Квантовая механика 1. Решения ДЗ №#1}}%
  \begin{center}
    {\small Квантовая механика 1 \enspace\textperiodcentered\enspace Семестр 1\par}
    \vspace{1ex}
    {\LARGE\bfseries Решения домашнего задания №#1\par}
    \vspace{0.5ex}
    {\large #2\par}
  \end{center}
  \noindent{\small Условия приведены в файле «#3», #4.}\par\medskip
}

\begin{document}
\solutiontitle{1}{Линейные операторы и коммутаторы}{ДЗ2026(1)\_01.pdf}{с.~1}

\section{Две эрмитовы матрицы}

\textbf{(a)} У \(K_1\) вещественные симметричные элементы. У \(K_2\)
выполнено \((K_2)_{ji}=(K_2)_{ij}^*\), поэтому обе матрицы равны своим
эрмитовым сопряжениям. Перемножение в двух порядках даёт
\[
 K_1K_2=\begin{pmatrix}i&0&-i\\0&0&0\\i&0&-i\end{pmatrix},
 \qquad K_2K_1=\begin{pmatrix}-i&0&-i\\0&0&0\\i&0&i\end{pmatrix},
 \qquad [K_1,K_2]=2i\,\operatorname{diag}(1,0,-1).
\]

\textbf{(b)} Из \(K_1(u,v,w)^{\mathsf T}=(v,u+w,v)^{\mathsf T}\) следуют
уравнения \(v=\lambda u=\lambda w\), \(u+w=\lambda v\).
При \(\lambda=0\) они дают \(v=0,\ w=-u\). При \(\lambda\ne0\)
имеем \(w=u\) и \(\lambda^2=2\). После деления на длины векторов получаем
\[
 \lambda_0=0,\quad v_0=\frac{(1,0,-1)^{\mathsf T}}{\sqrt2},
 \qquad \lambda_\pm=\pm\sqrt2,\quad
 v_\pm=\frac{(1,\pm\sqrt2,1)^{\mathsf T}}{2}.
\]
Три вектора с различными собственными значениями образуют базис.

\section{Тождества для коммутаторов}

Для неограниченных операторов все произведения в этом пункте
понимаются на общей инвариантной области.

\textbf{(a)} \( [A,B]=AB-BA=-(BA-AB)=-[B,A]\).

\textbf{(b)} По дистрибутивности
\begin{align*}
 [A+B,C]&=AC+BC-CA-CB=[A,C]+[B,C],\\
 [A,B+C]&=AB+AC-BA-CA=[A,B]+[A,C].
\end{align*}

\textbf{(c)} В первом коммутаторе добавляем и вычитаем \(BAC\),
во втором проделываем то же с \(ACB\). Получаем
\begin{align*}
 [A,BC]&=(ABC-BAC)+(BAC-BCA)=[A,B]C+B[A,C],\\
 [AB,C]&=(ABC-ACB)+(ACB-CAB)=A[B,C]+[A,C]B.
\end{align*}

\textbf{(d)} Применяя обе формулы пункта (c) в разном порядке, получаем
\begin{align*}
 [AB,CD]&=A[B,CD]+[A,CD]B\\
        &=A[B,C]D+AC[B,D]+[A,C]DB+C[A,D]B,\\
 [AB,CD]&=[AB,C]D+C[AB,D]\\
        &=A[B,C]D+[A,C]BD+CA[B,D]+C[A,D]B .
\end{align*}
В каждом раскрытии порядок множителей сохранён.

\section{Дифференциальные операторы и их спектры}

\textbf{(a)} При \(A^2v=\lambda v\), \(0\ne v\in D(A^2)\), симметричность
даёт \(\lambda\|v\|^2=\braket{v}{A^2v}=\braket{Av}{Av}\ge0\).
Отсюда \(\lambda\ge0\).

\textbf{(b)} Возьмём стандартное пространство \(L^2(\mathbb R,\dd x)\).
Одной ограниченности функций из условия недостаточно для существования
скалярного произведения. Например, \(1\) имеет бесконечную норму.
На пространстве Шварца граничные члены исчезают, и интегрирование по
частям даёт
\begin{align*}
 \int f^*(-ig')\,\dd x
   &=[-if^*g]_{-\infty}^{\infty}+\int ( -if')^*g\,\dd x,\\
 \int f^*(-g'')\,\dd x
   &=[-f^*g'+f'^*g]_{-\infty}^{\infty}+\int(-f'')^*g\,\dd x.
\end{align*}
Для самосопряжённости берём естественные области \(H^1(\mathbb R)\) и
\(H^2(\mathbb R)\) соответственно. Унитарное преобразование Фурье
переводит эти операторы в умножение на вещественные \(k\) и \(k^2\)
на максимальных областях \(k\tilde f\in L^2\) и \(k^2\tilde f\in L^2\).
Оператор умножения на вещественную функцию самосопряжён на такой области.

\textbf{(c)} На пробных функциях сопряжённое дифференциальное выражение
для \(F=\alpha x+i\beta\partial_x\) имеет коэффициенты
\(\alpha^*,\beta^*\). Выражения \(x\) и \(i\partial_x\) линейно
независимы, поэтому условие эрмитовости требует
\(\alpha,\beta\in\mathbb R\).
При \(\beta\ne0\) решение \(i\beta f'=(\lambda-\alpha x)f\) имеет вид
\[
 f_\lambda(x)=C_\lambda
 \exp\!\left(\frac{i\alpha x^2}{2\beta}-\frac{i\lambda x}{\beta}\right).
\]
Для вещественных коэффициентов умножение \(U\) на фазу
\(e^{i\alpha x^2/(2\beta)}\) унитарно. По правилу производной
\(F(Ug)=U(i\beta g')\), так что самосопряжённая реализация
\(F=U(i\beta\partial_x)U^\dagger\) имеет область \(UH^1\)
и непрерывный спектр \(\mathbb R\). Выбор \(C_\lambda=(2\pi|\beta|)^{-1/2}\)
даёт ортогональность и полноту в смысле распределений.
\begin{align*}
 \int f_\lambda(x)^*f_\mu(x)\,\dd x
  &=\frac{1}{2\pi|\beta|}\int e^{i(\lambda-\mu)x/\beta}\,\dd x
    =\delta(\lambda-\mu),\\
 \int f_\lambda(x)^* f_\lambda(x')\,\dd\lambda
  &=e^{i\alpha(x'^2-x^2)/(2\beta)}\delta(x-x')=\delta(x-x').
\end{align*}
В последнем равенстве фазовый множитель равен единице на носителе
\(\delta(x-x')\).

При \(\beta=0,\ \alpha\ne0\) оператор умножает на \(\alpha x\).
Базис \(f_\xi(x)=\delta(x-\xi)\) пометим координатой \(\xi\),
для которой \(\lambda=\alpha\xi\).
Интегрирование по \(x\) даёт \(\delta(\xi-\xi')\), а интегрирование по
\(\xi\) даёт \(\delta(x-x')\). Для вещественного \(\alpha\) это то же
разложение по спектральному параметру \(\lambda\) с
\(f_\lambda(x)=|\alpha|^{-1/2}\delta(x-\lambda/\alpha)\).
При \(\alpha=\beta=0\) весь \(L^2\) является собственным пространством
для \(\lambda=0\), и можно выбрать любой ортонормированный базис.

Если допускаются комплексные коэффициенты при \(\beta\ne0\), положим
\(c=\operatorname{Im}(\alpha/\beta)\). При \(c=0\) приведённое
унитарное преобразование по-прежнему работает со спектральной меткой
\(s=\lambda/\beta\in\mathbb R\) и мерой \(\dd s\).
При \(c>0\) модуль решения пропорционален
\(\exp(-cx^2/2+\operatorname{Im}(\lambda/\beta)x)\).
Поэтому на максимальной области \(f\in L^2,\ Ff\in L^2\) каждая
\(\lambda\in\mathbb C\) является собственным значением.
Скалярное произведение любых двух таких решений есть ненулевой
гауссов интеграл \(\int e^{-cx^2+bx}\dd x=\sqrt{\pi/c}\,e^{b^2/(4c)}\).
Семейство имеет плотную линейную оболочку. Действительно,
ортогональность всем \(f_{\beta s}\) с \(s\in\mathbb R\) обнуляет
преобразование Фурье произведения проверяемой функции на сопряжённый
гауссов множитель, и этот множитель нигде не равен нулю.
При \(c<0\) решения растут хотя бы на одном конце оси и в \(L^2\)
собственных векторов нет. Требуемое ортонормированное разложение по
вещественным \(\lambda\) относится к эрмитову случаю.

\section{Унитарные операторы}

\textbf{(a)} Если \(Uv=\lambda v\) и \(v\ne0\), то из \(U^\dagger U=I\)
следует \(\|v\|^2=\|Uv\|^2=|\lambda|^2\|v\|^2\), откуда
\(|\lambda|=1\).

\textbf{(b*)} Для матриц или ограниченного \(A=A^\dagger\) сопряжение
степенного ряда даёт \((e^{iA})^\dagger=e^{-iA}\). При перемножении
этих рядов коэффициент у \(A^n\) равен
\(i^n\sum_{k=0}^n\binom nk(-1)^k/n!\), то есть \(1\) при \(n=0\)
и \(0\) при \(n>0\). Получаем \(e^{-iA}e^{iA}=e^{iA}e^{-iA}=I\).
Для неограниченного самосопряжённого \(A\) тот же результат следует
из спектрального исчисления, поскольку \(e^{i\lambda}\) имеет модуль
один при всех вещественных \(\lambda\).

\end{document}
